\documentclass[a4paper,12pt,twoside]{memoir}

% Castellano
\usepackage[spanish,es-tabla]{babel}
\selectlanguage{spanish}
\usepackage[utf8]{inputenc}
\usepackage[T1]{fontenc}
\usepackage{lmodern} % Scalable font
\usepackage{microtype}
\usepackage{placeins}

\RequirePackage{booktabs}
\RequirePackage[table]{xcolor}
\RequirePackage{xtab}
\RequirePackage{multirow}

% Links
\PassOptionsToPackage{hyphens}{url}\usepackage[colorlinks]{hyperref}
\hypersetup{
	allcolors = {red}
}

% Ecuaciones
\usepackage{amsmath}

% Rutas de fichero / paquete
\newcommand{\ruta}[1]{{\sffamily #1}}

% Párrafos
\nonzeroparskip

% Huérfanas y viudas
\widowpenalty100000
\clubpenalty100000

% Imágenes

% Comando para insertar una imagen en un lugar concreto.
% Los parámetros son:
% 1 --> Ruta absoluta/relativa de la figura
% 2 --> Texto a pie de figura
% 3 --> Tamaño en tanto por uno relativo al ancho de página
\usepackage{graphicx}
\newcommand{\imagen}[3]{
	\begin{figure}[!h]
		\centering
		\includegraphics[width=#3\textwidth]{#1}
		\caption{#2}\label{fig:#1}
	\end{figure}
	\FloatBarrier
}

% Comando para insertar una imagen sin posición.
% Los parámetros son:
% 1 --> Ruta absoluta/relativa de la figura
% 2 --> Texto a pie de figura
% 3 --> Tamaño en tanto por uno relativo al ancho de página
\newcommand{\imagenflotante}[3]{
	\begin{figure}
		\centering
		\includegraphics[width=#3\textwidth]{#1}
		\caption{#2}\label{fig:#1}
	\end{figure}
}

% El comando \figura nos permite insertar figuras comodamente, y utilizando
% siempre el mismo formato. Los parametros son:
% 1 --> Porcentaje del ancho de página que ocupará la figura (de 0 a 1)
% 2 --> Fichero de la imagen
% 3 --> Texto a pie de imagen
% 4 --> Etiqueta (label) para referencias
% 5 --> Opciones que queramos pasarle al \includegraphics
% 6 --> Opciones de posicionamiento a pasarle a \begin{figure}
\newcommand{\figuraConPosicion}[6]{%
  \setlength{\anchoFloat}{#1\textwidth}%
  \addtolength{\anchoFloat}{-4\fboxsep}%
  \setlength{\anchoFigura}{\anchoFloat}%
  \begin{figure}[#6]
    \begin{center}%
      \Ovalbox{%
        \begin{minipage}{\anchoFloat}%
          \begin{center}%
            \includegraphics[width=\anchoFigura,#5]{#2}%
            \caption{#3}%
            \label{#4}%
          \end{center}%
        \end{minipage}
      }%
    \end{center}%
  \end{figure}%
}

%
% Comando para incluir imágenes en formato apaisado (sin marco).
\newcommand{\figuraApaisadaSinMarco}[5]{%
  \begin{figure}%
    \begin{center}%
    \includegraphics[angle=90,height=#1\textheight,#5]{#2}%
    \caption{#3}%
    \label{#4}%
    \end{center}%
  \end{figure}%
}
% Para las tablas
\newcommand{\otoprule}{\midrule [\heavyrulewidth]}
%
% Nuevo comando para tablas pequeñas (menos de una página).
\newcommand{\tablaSmall}[5]{%
 \begin{table}
  \begin{center}
   \rowcolors {2}{gray!35}{}
   \begin{tabular}{#2}
    \toprule
    #4
    \otoprule
    #5
    \bottomrule
   \end{tabular}
   \caption{#1}
   \label{tabla:#3}
  \end{center}
 \end{table}
}

%
% Nuevo comando para tablas pequeñas (menos de una página).
\newcommand{\tablaSmallSinColores}[5]{%
 \begin{table}[H]
  \begin{center}
   \begin{tabular}{#2}
    \toprule
    #4
    \otoprule
    #5
    \bottomrule
   \end{tabular}
   \caption{#1}
   \label{tabla:#3}
  \end{center}
 \end{table}
}

\newcommand{\tablaApaisadaSmall}[5]{%
\begin{landscape}
  \begin{table}
   \begin{center}
    \rowcolors {2}{gray!35}{}
    \begin{tabular}{#2}
     \toprule
     #4
     \otoprule
     #5
     \bottomrule
    \end{tabular}
    \caption{#1}
    \label{tabla:#3}
   \end{center}
  \end{table}
\end{landscape}
}

%
% Nuevo comando para tablas grandes con cabecera y filas alternas coloreadas en gris.
\newcommand{\tabla}[6]{%
  \begin{center}
    \tablefirsthead{
      \toprule
      #5
      \otoprule
    }
    \tablehead{
      \multicolumn{#3}{l}{\small\sl continúa desde la página anterior}\\
      \toprule
      #5
      \otoprule
    }
    \tabletail{
      \hline
      \multicolumn{#3}{r}{\small\sl continúa en la página siguiente}\\
    }
    \tablelasttail{
      \hline
    }
    \bottomcaption{#1}
    \rowcolors {2}{gray!35}{}
    \begin{xtabular}{#2}
      #6
      \bottomrule
    \end{xtabular}
    \label{tabla:#4}
  \end{center}
}

%
% Nuevo comando para tablas grandes con cabecera.
\newcommand{\tablaSinColores}[6]{%
  \begin{center}
    \tablefirsthead{
      \toprule
      #5
      \otoprule
    }
    \tablehead{
      \multicolumn{#3}{l}{\small\sl continúa desde la página anterior}\\
      \toprule
      #5
      \otoprule
    }
    \tabletail{
      \hline
      \multicolumn{#3}{r}{\small\sl continúa en la página siguiente}\\
    }
    \tablelasttail{
      \hline
    }
    \bottomcaption{#1}
    \begin{xtabular}{#2}
      #6
      \bottomrule
    \end{xtabular}
    \label{tabla:#4}
  \end{center}
}

%
% Nuevo comando para tablas grandes sin cabecera.
\newcommand{\tablaSinCabecera}[5]{%
  \begin{center}
    \tablefirsthead{
      \toprule
    }
    \tablehead{
      \multicolumn{#3}{l}{\small\sl continúa desde la página anterior}\\
      \hline
    }
    \tabletail{
      \hline
      \multicolumn{#3}{r}{\small\sl continúa en la página siguiente}\\
    }
    \tablelasttail{
      \hline
    }
    \bottomcaption{#1}
  \begin{xtabular}{#2}
    #5
   \bottomrule
  \end{xtabular}
  \label{tabla:#4}
  \end{center}
}



\definecolor{cgoLight}{HTML}{EEEEEE}
\definecolor{cgoExtralight}{HTML}{FFFFFF}

%
% Nuevo comando para tablas grandes sin cabecera.
\newcommand{\tablaSinCabeceraConBandas}[5]{%
  \begin{center}
    \tablefirsthead{
      \toprule
    }
    \tablehead{
      \multicolumn{#3}{l}{\small\sl continúa desde la página anterior}\\
      \hline
    }
    \tabletail{
      \hline
      \multicolumn{#3}{r}{\small\sl continúa en la página siguiente}\\
    }
    \tablelasttail{
      \hline
    }
    \bottomcaption{#1}
    \rowcolors[]{1}{cgoExtralight}{cgoLight}

  \begin{xtabular}{#2}
    #5
   \bottomrule
  \end{xtabular}
  \label{tabla:#4}
  \end{center}
}



\graphicspath{ {./img/} }

% Capítulos
\chapterstyle{bianchi}
\newcommand{\capitulo}[2]{
	\setcounter{chapter}{#1}
	\setcounter{section}{0}
	\setcounter{figure}{0}
	\setcounter{table}{0}
	\chapter*{\thechapter.\enskip #2}
	\addcontentsline{toc}{chapter}{\thechapter.\enskip #2}
	\markboth{#2}{#2}
}

% Apéndices
\renewcommand{\appendixname}{Apéndice}
\renewcommand*\cftappendixname{\appendixname}

\newcommand{\apendice}[1]{
	%\renewcommand{\thechapter}{A}
	\chapter{#1}
}

\renewcommand*\cftappendixname{\appendixname\ }

% Formato de portada
\makeatletter
\usepackage{xcolor}
\newcommand{\tutor}[1]{\def\@tutor{#1}}
\newcommand{\course}[1]{\def\@course{#1}}
\definecolor{cpardoBox}{HTML}{E6E6FF}
\def\maketitle{
  \null
  \thispagestyle{empty}
  % Cabecera ----------------
\noindent\includegraphics[width=\textwidth]{cabecera}\vspace{1cm}%
  \vfill
  % Título proyecto y escudo informática ----------------
  \colorbox{cpardoBox}{%
    \begin{minipage}{.8\textwidth}
      \vspace{.5cm}\Large
      \begin{center}
      \textbf{TFG del Grado en Ingeniería Informática}\vspace{.6cm}\\
      \textbf{\LARGE\@title{}}
      \end{center}
      \vspace{.2cm}
    \end{minipage}

  }%
  \hfill\begin{minipage}{.20\textwidth}
    \includegraphics[width=\textwidth]{escudoInfor}
  \end{minipage}
  \vfill
  % Datos de alumno, curso y tutores ------------------
  \begin{center}%
  {%
    \noindent\LARGE
    Presentado por \@author{}\\ 
    en Universidad de Burgos --- \@date{}\\
    Tutor: \@tutor{}\\
  }%
  \end{center}%
  \null
  \cleardoublepage
  }
\makeatother

\newcommand{\nombre}{Daniel Fernández Barrientos} %%% cambio de comando
\newcommand{\dni}{51079809-Y} %%% cambio de comando
\newcommand{\tfgTitulo}{Generador de cuestionarios y escenarios de test sobre soft skills} %%% cambio de comando
\newcommand{\tfgTutor}{Raúl Marticorena Sanchez} %%% cambio de comando

% Datos de portada
\title{\tfgTitulo}
\author{\nombre}
\tutor{\tfgTutor}
\date{\today}

\begin{document}

\maketitle


\newpage\null\thispagestyle{empty}\newpage


%%%%%%%%%%%%%%%%%%%%%%%%%%%%%%%%%%%%%%%%%%%%%%%%%%%%%%%%%%%%%%%%%%%%%%%%%%%%%%%%%%%%%%%%
\thispagestyle{empty}


\noindent\includegraphics[width=\textwidth]{cabecera}\vspace{1cm}

\noindent D. \tfgTutor, profesor del departamento de nombre departamento, área de nombre área.

\noindent Expone:

\noindent Que el alumno D. \nombre, con DNI \dni, ha realizado el Trabajo final de Grado en Ingeniería Informática titulado \tfgTitulo de TFG. 

\noindent Y que dicho trabajo ha sido realizado por el alumno bajo la dirección del que suscribe, en virtud de lo cual se autoriza su presentación y defensa.

\begin{center} %\large
En Burgos, {\large \today}
\end{center}

\vfill\vfill\vfill

% Author and supervisor
\begin{minipage}{0.45\textwidth}
\begin{flushleft} %\large
Vº. Bº. del Tutor:\\[2cm]
D. nombre tutor
\end{flushleft}
\end{minipage}
\hfill
\begin{minipage}{0.45\textwidth}
\begin{flushleft} %\large
Vº. Bº. del co-tutor:\\[2cm]
D. nombre co-tutor
\end{flushleft}
\end{minipage}
\hfill

\vfill

% para casos con solo un tutor comentar lo anterior
% y descomentar lo siguiente
%Vº. Bº. del Tutor:\\[2cm]
%D. nombre tutor


\newpage\null\thispagestyle{empty}\newpage




\frontmatter

% Abstract en castellano
\renewcommand*\abstractname{Resumen}
\begin{abstract}
El proyecto consiste en la creación de una aplicación web que permita la realización de cuestionarios, donde la respuesta de cada pregunta tiene asociada una puntuación.
La aplicación guarda registros de todos los cuestionarios realizados y permite exportarlos para su explotación.
\end{abstract}

\renewcommand*\abstractname{Descriptores}
\begin{abstract}
Spring boot, java, docker, eclipse, maven, latex, mysql, github, zube, json, csv, html, css, bootstrap, thymeleaf,  spring jpa, spring web, spring security, email
\end{abstract}

\clearpage

% Abstract en inglés
\renewcommand*\abstractname{Abstract}
\begin{abstract}
The project consists of the creation of a web application that allows questionnaires to be carried out, where each question's answer has a score associated with it.
The application stores all the information about the questionnaires that have been carried out and allows its export.
\end{abstract}

\renewcommand*\abstractname{Keywords}
\begin{abstract}
Spring boot, java, docker, eclipse, maven, latex, mysql, github, zube, json, csv, html, css, bootstrap, thymeleaf,  spring jpa, spring web, spring security, email
\end{abstract}

\clearpage

% Indices
\tableofcontents

\clearpage

\listoffigures

\clearpage

\listoftables
\clearpage

\mainmatter
\capitulo{1}{Introducción}

En este primer apartado se va a tratar sobre la descripción del contenido del trabajo, de la estructura de la memoria y del resto de materiales entregados.

\capitulo{2}{Objetivos del proyecto}

Este apartado explica de forma precisa y concisa cuales son los objetivos que se persiguen con la realización del proyecto. Se puede distinguir entre los objetivos marcados por los requisitos del software a construir y los objetivos de carácter técnico que plantea a la hora de llevar a la práctica el proyecto.

\capitulo{3}{Conceptos teóricos}

En aquellos proyectos que necesiten para su comprensión y desarrollo de unos conceptos teóricos de una determinada materia o de un determinado dominio de conocimiento, debe existir un apartado que sintetice dichos conceptos.

Algunos conceptos teóricos de \LaTeX{} \footnote{Créditos a los proyectos de Álvaro López Cantero: Configurador de Presupuestos y Roberto Izquierdo Amo: PLQuiz}.

\section{Secciones}

Las secciones se incluyen con el comando section.

\subsection{Subsecciones}

Además de secciones tenemos subsecciones.

\subsubsection{Subsubsecciones}

Y subsecciones. 


\section{Referencias}

Las referencias se incluyen en el texto usando cite~\cite{udemy:eliseo}. Para citar webs, artículos o libros~\cite{koza92}, si se desean citar más de uno en el mismo lugar~\cite{bortolot2005, bortolot2008}.


\section{Imágenes}

Se pueden incluir imágenes con los comandos standard de \LaTeX, pero esta plantilla dispone de comandos propios como por ejemplo el siguiente:

\imagen{escudoInfor}{Autómata para una expresión vacía}{.5}



\section{Listas de items}

Existen tres posibilidades:

\begin{itemize}
	\item primer item.
	\item segundo item.
\end{itemize}

\begin{enumerate}
	\item primer item.
	\item segundo item.
\end{enumerate}

\begin{description}
	\item[Primer item] más información sobre el primer item.
	\item[Segundo item] más información sobre el segundo item.
\end{description}
	
\begin{itemize}
\item 
\end{itemize}

\section{Tablas}

Igualmente se pueden usar los comandos específicos de \LaTeX o bien usar alguno de los comandos de la plantilla.

\tablaSmall{Herramientas y tecnologías utilizadas en cada parte del proyecto}{l c c c c}{herramientasportipodeuso}
{ \multicolumn{1}{l}{Herramientas} & App AngularJS & API REST & BD & Memoria \\}{ 
HTML5 & X & & &\\
CSS3 & X & & &\\
BOOTSTRAP & X & & &\\
JavaScript & X & & &\\
AngularJS & X & & &\\
Bower & X & & &\\
PHP & & X & &\\
Karma + Jasmine & X & & &\\
Slim framework & & X & &\\
Idiorm & & X & &\\
Composer & & X & &\\
JSON & X & X & &\\
PhpStorm & X & X & &\\
MySQL & & & X &\\
PhpMyAdmin & & & X &\\
Git + BitBucket & X & X & X & X\\
Mik\TeX{} & & & & X\\
\TeX{}Maker & & & & X\\
Astah & & & & X\\
Balsamiq Mockups & X & & &\\
VersionOne & X & X & X & X\\
} 

\capitulo{4}{Técnicas y herramientas}


En esta parte de la memoria se van a presentar las distintas herramientas y técnicas utilizadas en el desarollo del proyecto.
Se incluirán todas las que se hayan valorado o utilizado, incluso aquellas que tras ciertas pruebas o pasos, se hayan descartado por haber encontrado una alternativa más funcional.

\section{Spring Tool Suite for Eclipse}

Herramienta de desarrollo (IDE) que se utilizará para llevar a cambo el proyecto.
A pesar de haber otras similares, como InteliJ, he preferido utilizar la \textit{suite} de \textit{Spring Boot} para Eclipse al estar más fmailiarizado con el entorno de trabajo, ya que ha sido la herramienta utilizada en otras asignaturas del grado.

\section{\LaTeX}
El editor de texto para generar toda la documentación relacionada con el trabajo de fin de grado será \LaTeX.
Aún estándo acostumbrado a utilizar Microsoft Word para este tipo de tareas, veo en utilizar \LaTeX  una oportunidad de aprender un nuevo generador de documentos.
La curva de aprendizaje es más grande que con Word (en mi caso), pero el resultado merece la pena porque permite aplicar formatos de forma sencilla, un aspecto en el que Word para documentos muy grandes, está en clara desventaja.

\section{Mik\TeX{}}
Para trabajar con plantillas de \LaTeX  en entornos windows existen varias alternativas, en mi caso he preferido utilizar el programa Mik\TeX{}.
Es un programa sencillo y tiene por separado tanto la parte de edición como la de presentación, para ir viendo los cambios una vez se compila la plantilla completa.

\section{Spring Boot}

Para la aplicación web en java se utilizará el Framework de Spring.
Es mi primera vez con un framework en java, y para aprender conceptos de forma más rápida y con ejemplos he tomado como referencia un curso en Udemy \cite{udemy:eliseo} para dar un acelerón en el aprendizaje.
Gracias a las dependencias predefinidas, facilita el desarrollo de la aplicación de forma considerable.

\subsection{Spring Web}

Es la dependencia principal de Spring Boot para la creación de aplicaciones web. 
Simplifica la creación de controladores, el manejo de las solicitude y se integra de forma sencilla con motores de plantillas como Thymeleaf.

\subsection{Spring JPA}

Es la dependencia de Spring para el manejo de entidades persistentes de una base de datos. 
Con los respositorios CRUD o JPA, y la inclusión por defecto de Hibernate, simplifica hasta el extremo la gestión de nuestros los objetos en la base de datos.

\subsection{Spring Security}

Es la dependencia de Spring para dotar de seguridad nuestra aplicación de forma automática.

Esta seguridad se basa en autenticación y autorizació principalmente, para poder permitir ciertas partes de nuestra web a usuarios "invitados" o limitar el acceso a otras en función del rol del usuario que haya iniciado sesión.

Integra de forma predeterminada protección contra amenazas comunes, como CSRF (Cross-Site Request Forgery).

\section{Thymeleaf}

Es un motor de plantillas para HTML, para permitir la creación de código HTML de forma dinámica.
Es un motor sencillo de usar, con expresiones básicas, que nos permite utilizar los objetos que enviamos para generar más o menos código, mostrar u ocultar opciones, \dots
Su curva de aprendizaje no es sencilla, pero dedicándole un poco de tiempo, se ahorra mucho a la hora de generar el código de las páginas web.
El uso de los \textit{fragments} es un gran aliado para ahorrarnos líneas de código.

\section{HTML}

Lenguaje utilizado en la creación de las páginas web, en particular HTML5, que tiene algunas funcionalidades básicas ya implementadas, como por ejemplo el uso de "required" en campos de formulario para que sea obligatorio rellenar los campos si necesidad de utilizar javascript para el control.

\section{MySQL}

Se utiliza MySQL como base de datos para almacenar los usuarios y roles.
Se ha decidido usar una base de datos porque su integración con Spring Security es muy sencilla y funcional. 
A través de los repositorios CRUD que nos proporciona Spring JPA podemos gestionar las tablas y registros de forma ágil.

\section{Docker}

\textbf{Docker} consiste en tener microservicios corriendo en una máquina, si necesidad de tener una gran infraestructura montada.

He utilizado docker tanto para el servidor de MySQL como para la propia aplicación Web, de esta forma con 2 simples docker, un docker-compose (fichero de configuración de dockers que permite descargarse, crear, parar o lanzar dockers) podemos tener toda la aplicación completa corriendo en cualquier entorno, ya que docker funciona en los principales sistemas operativos del mercado: Linux, Windows y macOS.

\subsection{Docker Compose}

Esta herramienta complementaria de docker, \textit{docker compose} \cite{web:dockercompose} permite definir y gestionar aplicaciones que requieran más de un contenedor, de forma conjunta y sencilla.
Se ha utilizado como base para la instalación del proyecto, explicando sus funcionalidades y permitiendo ciertas personalizaciones en el despligue de la aplicación.

\section{GitHub}

Se va utiliza GitHub para dar visibilidad al trabajo diario y constante en el proyecto, ya que permite realizar aportaciones incrementales de código, documentación, etc del TFG completo.

Es la herramienta por excelencia para llevar un control de versiones en cualquier proyecto, sea software o no.

\subsection{Repository}

Se ha creado un repositorio~\cite{github:repo} donde se subirá todo el TFG al completo, tanto la documentación de la memoria, como la aplicación de java.

\subsection{Releases}

En esta sección dentro del repositorio se han ido subiendo las diferentes versiones funcionales de la aplicación, para tener un lugar central desde el que descargarse la última versión ``liberada''.

\subsection{Packages}

En esta parte de GitHub se han subido las imágenes de docker funcionales de la aplicación, para pemitir un despliegue sencillo desde una ubicación centralizada.
Con un simple comando de \textit{docker pull url de la imagen} puede descargarte la imagen para empezar a utilizarla.

\subsection{Proyect}

Se ha creado un proyecto~\cite{github:proyect} relacionado, que sigue como la base la metodología Kanban, para crear nuevas tareas, relacionarlas con el repositorio y utilizar pizarras y paneles para ver el progreso, en lo que se está trabajando y lo ya completado.

\textbf{Se prescinde de esta herramienta al sustituirse por Zube.io}

\section{JSON}

Se van a utilizar ficheros con la estructura de un json para almacenar los cuestionarios completos, a los que se accederá desde la aplicación web de Java.

\section{Zube.io}

Se va a utilizar la versión gratuita de Zube.io, que permite generar Sprints, y gráficos de seguimiento, como \textit{Burndown} o \textit{Burnup}.

Gracias a la sencilla y completa integración con Github ofrece una versión más completa de la parte de Proyect de Github, por lo que he decidido utilizar esta plataforma para el seguimiento de las tareas en vez de la que ofrecía Github.

La versión gratuita permite integrar lo necesario del repositorio y del proyecto, por lo que no es necesario ampliar a versiones más completas.

Como pega, no permite hacerlo público para compartir su contenido, pero con los gráficos se suple esta limitación.

\section{GitPod}

Se ha estudiado la posibilidad de usar \textit{gitpod} para alojar una versión funcional de la aplicación, ya que diferentes alternativas como \textit{Heroku} o \textit{Render} ni permitían un fácil despliegue con docker-compose.yml ni daban la opción de una prueba real de forma gratuita.

A pesar de haber vinculado \textit{gitpod} con \textit{GitHub}, se limitaba a 50 horas el despliegue de una aplicación, pero eso no se incluye como parte del proyecto.









\capitulo{5}{Aspectos relevantes del desarrollo del proyecto}

Este apartado pretende recoger los aspectos más interesantes del desarrollo del proyecto, incluyendo las diferentes decisiones que se han ido tomando y los problemas que se han tenido que afrontar.

\section{Inicio del proyecto}

La selección del proyecto suponía un doble reto personal:

\begin{itemize}
	\item Nunca había programado aplicaciones web ni conocía la tecnología de Spring Boot.
	\item No había un proyecto base, documentado, como inicio del proyecto, sino que era crear algo nuevo desde cero.
\end{itemize}

Tras la primera reunión con el tutor, donde me expuso el camino que se quería seguir con el proyecto, me puse a recabar información sobre ello.

\section{Metodologías}

Desde el inicio del proyecto se ha dedicado tiempo y esfuerzo en que la realización del proyecto fuese de la manera más profesional posible.

La gestión del proyecto se ha realizado teniendo como base la metodología ágil de \textit{Scrum} ~\cite{web:scrum}.

Al tratarse de un Trabajo de Fin de Grado, que se realiza de forma individual, no se siguió de forma completa, porque el grupo de desarrolladores estaba formado solo por un miembro.

Cómo se ha aplicado la metodología:

\begin{itemize}
	\item Reunión inicial de \textit{kick-off} ~\cite{web:kickoff} del proyecto, donde se presentó el proyecto y se establecieron ciertos objetivos generales.
	\item Se utilizó una estrategia de desarrollo incremental basada en \textit{sprints}, a través de iteraciones y revisiones.
	\item La duración media de los \textit{sprints} fue de dos semanas, excepto los últimos del proyecto, que han sido de una semana.
	\item En la finalización de cada \textit{sprints} se entregaba un incremento del proyecto.
	\item Se realizaba una reunión de revisión del \textit{sprint} que servía también como reunión de planificación del nuevo \textit{sprint}.
	\item Se definía en el tablero kanban de \textit{Zube.io} una lista de tareas y se asociaban al nuevo \textit{sprint}.
	\item Para monitorizar el progreso se han utilizado gráficos \textit{burnup}.
\end{itemize}

\section{Formación}

La principal tecnología en el desarrollo del proyecto ha sido \textit{Spring Boot}, sobre la que no se tenía ningún conocimiento previo, lo que me ha llevado a tener que formarme por diferentes medios:

\begin{itemize}
	\item Spring Boot 3. Aplicaciones web y REST APIs con Spring MVC (Udemy) ~\cite{udemy:eliseo}.
	\item Cómo convertir en imagen de docker una aplicación de Spring Boot ~\cite{doc:springbootdocker}.
	\item Cómo aplicar estilos de personalización de \textit{bootstrap} y \textit{CSS} en documentos \textit{HTML}.
\end{itemize}

Me gustaría destacar la web \href{https://www.baeldung.com/spring-boot}{Baeldung}, al contener mucha información sobre Spring Boot en general.

Aunque en menor menida, también he tenido que formarme en otras de las tecnologías utilizadas:

\begin{itemize}
	\item Uso y configuración de un archivo \textit{docker compose} ~\cite{web:dockercompose}.
	\item Cómo se utiliza \LaTeX para la generación de documentos.
\end{itemize}

\section{Desarrollo de la aplicación web}

Se va a dividir el desarrollo de la aplicación web en varios apartados, para poder detallar mejor el proceso de construcción de la misma.

\subsection{Código de la aplicación}

El desarrollo del código de la aplicación supuso un primer reto, pues en el inicio del proyecto todavía no terminaba de comprender cómo iba a crear una aplicación web con Java, ya que no era lo que estaba acostumbrado a crear con este lenguaje de programación.

Una vez se comprendió cómo funcionaba, de forma básica, \textit{Spring Web} como \textit{starter} del proyecto, este paso se convirtió en algo trivial.

La primera parte de la creación de la aplicación consistía en poder ``leer'' un archivo \textit{JSON} y convertirlo a objetos en Java con los que se pudiese trabajar de forma nativa.

Se crearon varias clases modelo para agrupar el contenido del \textit{JSON}:

 \begin{itemize}
	\item Clase \textit{Quiz}: Contiene la descripción básica del cuestionario y un listado de preguntas.
	\item Clase \textit{Questions}: Contiene la información de cada pregunta y un listado de respuestas.
	\item Clase \textit{Answers}: Contiene la información de cada respuesta.
\end{itemize}

Una vez conseguido lo anterior, la base del proyecto ya estaba construida.

Para la presentación de contenidos en las páginas web, se usaron clases \textit{controller} en Java, que sirven de enlace entre la propia aplicación y el código \textit{HTML} de la página web.

A través de estas clases, se pasaban los cuestionarios completos al código \textit{HTML} de la página web y en ella, con \textit{Thymeleaf} se mostraba la información.

Ya se tenía la forma de ver el contenido del \textit{JSON} en las páginas web.

Hasta aquí tendríamos una base más ``robusta'' de la aplicación, teniendo todo lo necesario dentro del propio proyecto de Java.

Los primeros problemas los encontramos al intentar actualizar los archivos de propiedades, ya que al ser internos de la aplicación, necesitaban un reinicio de la misma cada vez que había un cambio, algo que no era funcional para un entorno de producción.

A raíz de lo anterior, se empezó a trabajar con archivos \textit{JSON} externos a la aplicación, utilizando rutas relativas para poder encontrar los archivos siempre, independientemene de la ubicación raíz del proyecto.

Finalmente, se incluyeron varios tipos de clases de Java adicionales:

 \begin{itemize}
	\item Clases de configuración: Se utilizaron ciertas clases para modificar la configuración por defecto de Spring Boot. También se ha utilizado para acceder a un archivo de propiedades personalizado, donde se definen algunos elementos de configuración y personalización de la aplicación.
	\item Clases de servicio: Contienen los métodos a los que se tienen que acceder desde otras clases, evitando duplicidad en el código y teniendo una clase centralizada para la gestión de la mayoría de los procesos, siendo más sencillo comprender su código y ahorrando métodos adicionales en otras clases.
\end{itemize}

\subsection{Código de las páginas web}

La estética y el diseño de las páginas web ha sido lo más complejo en cuanto a presentación final del proyecto.

Ayudado de \textit{bootstrap} y personalización a través de plantillas \textit{CSS} se ha conseguido aplicar un estilo homogéneo a toda la aplicación.

El uso de los fragmentos de \textit{Thymeleaf} ha ayudado a que los ficheros \textit{HTML} que componen la página web sean más legibles, evitando tener que escribir las mismas líneas de forma repetitiva, simplificándolo todo a una simple línea que enlaza con el \textit{fragmento} indicado.

Definitivamente, implementar las etiquetas \textit{HTML} ha sido lo más costoso en cuanto a programación se refiere, porque conseguir cuadrar y encajar todos los apartados en su sitio ha sido un trabajo de prueba y error constante.

\subsection{Integración de servicios y dependencias}

Los servicios o dependencias más destacables que implementa la aplicación son los siguientes:

\subsubsection{Servicio de envío de correos electrónicos}

Durante el proceso de desarrollo de un servicio de correo electrónico que enviase a los usuarios una nueva contraseña en caso de que no recordasen la suya, se probaron configuraciones con varios servidores de correo, pero finalmente, por simplifidad y acceso a documentación, se utilizó \textit{Gmail} que permite la creación de cuentas gratuitas y el uso de sus servidores de correo desde cualquier \textit{API} siempre que se utilice un \textit{token} ~\cite{gmail:token}.

La configuración completa del servidor de correo, así como, sus credenciales, se independizaron de la aplicación en un archivo de propiedades, para permitir su modificación por cualquier otro servicio de correo, dando libertar al programador para cambiarlo sin necesidad de compilar de nuevo el proyecto.

\subsubsection{Seguridad con Spring Security y Spring JPA}

En primer lugar, tras aplicar \textit{Spring Security}, la mayoría de las páginas dejaron de funcionar, todas aquellas que eran de tipo \textit{POST}.

Tras investigar descubrí que esto sucedía porque la primera medida de seguridad que se aplica, es la necesidad de recibir un \textit{CSRF token} (\textit{Cross-Site Request Forgery}) ~\cite{web:csrf} para evitar ataques \textit{man in the middle} a los formularios de las aplicaciones web, al requerir un \textit{token} de autenticación, para comprobar la autenticidad de la comunicación.

Debido a esto, todos los formularios tienen configurado un campo ``oculto'' que envía este \textit{token} al controlador de la aplicación, para permitir el acceso.

Se decidió utilizar una base de datos \textit{mysql} que se levanta cada vez que se inicia la aplicación, almacena los datos en una carpeta externa a la misma, para tener persistencia, y que está vinculada con el la seguridad de la aplicación, ya que se definieron los campos de nuestras tablas de usuarios y perfiles que se iban a usar tanto para iniciar sesión, como para comprobar a qué parte de la aplicación tiene cada uno.

\section{Documentación}

Para generar la documentación del proyecto, se ha utilizado una plantilla de \LaTeX, que permite mantener una estructura uniforme y el mismo estilo durante la documentación de todo el proyecto, aportando un aura de profesionalidad y sin los poblemas y complicaciones de formatos que nos podemos encontrar en otros programas comerciales.

La documentación de la aplicación se ha generado desde el propio editor de código, \textit{Spring Tool Suite 4 for Eclipse}, siguiendo el formato \textit{javadoc} ~\cite{doc:javadoc}.

Este formato permite la visualización de todas las clases utilizadas en formato \textit{HTML}, pudiendo navegar por los diferentes métodos y obteniendo una breve descripción de cada uno, así como de los parmáetros de entrada que esperan o del retorno de cada método.
\clearpage
\section{Publicación}

Se ha intentado realizar la publicación de la aplicación en la plataforma en la nube de \textit{GitPod} ~\cite{doc:gitpod}, porque permitía la integración de un repositorio de \textit{GitHub} y poder desplegar tus aplicaciones con un archivo \textit{docker-compose.yml}, pero las limitaciones del plan gratuito en cuanto a horas mensuales de ejecución no han permitido tenerla disponible de forma continua, sí puediendo comprobarse su correcto despliegue y funcionamiento.

Despliegue de la aplicación con ``docker compose up -d'' en \textit{GitPod}:
\imagen{gitPodDespliegue}{Despliegue de la aplicación en la plataforma web de GitPod.}{.5}

Acceso a la aplicación web a través de la dirección web que genera \textit{GitPod}:
\imagen{gitPodWeb}{Funcionamiento de la aplicación web a través de GitPod.}{.5}





\capitulo{6}{Trabajos relacionados}

Este apartado sería parecido a un estado del arte de una tesis o tesina. En un trabajo final grado no parece obligada su presencia, aunque se puede dejar a juicio del tutor el incluir un pequeño resumen comentado de los trabajos y proyectos ya realizados en el campo del proyecto en curso. 

\capitulo{7}{Conclusiones y Líneas de trabajo futuras}

En este capítulo se exponen las conclusiones derivadas del trabajo del proyecto, así como las líneas de trabajo futuras por las que se puede dar continuidad al proyecto.

\section{Conclusiones}

Tras el desarrollo del proyecto se obtienen las siguientes conclusiones:

\begin{itemize}
	\item Se ha cumplido con el objetivo general del proyecto, teniendo ahora disponible una aplicación web que permite la realización y gestión de cuestionarios.
	\item Haber utilizado tecnologías como Spring Boot y Docker, le dan al proyecto un aura de modernidad, al ser tecnologías que están en expansión dentro del ámbito del desarrollo de aplicaciones.
	\item He podido aplicar conocimientos de gran parte de lo aprendido durante el grado, y aprender sobre muchos otros en los que no contaba con ninguna experiencia previa.
	\item Integrar la creación del proyecto con herramientas de control de versiones ha facilitado el desarrollo incremental del proyecto, así como un \textit{know-how} de buenas prácticas para cualquier proyecto \textit{software} que tenga que desarrollar en el futuro.
	\item Estimar el tiempo total dedicado a pruebas y prototipos de la aplicación es muy complicado, ya que en algunos momentos me he atascado con parte de la programación y he tenido que probar vías alternativas, pero el uso de la metodología ágil para la gestión de las tareas ha facilitado el seguimiento del progreso.
\end{itemize}

\section{Líneas de trabajo futuras}

La entrega del Trabajo de Fin de Grado solo es el inicio del proyecto, ya que su desarrollo contiúa.

A continuación, se resume el camino a seguir del proyecto:

\begin{itemize}
	\item Integrar en la página web la posibilidad de editar los cuestionarios ya cargados, sin tener que actualizar el archivo \textit{JSON} y volver a cargarlo.
	\item Integrar la aplicación con una \textit{API} de traducción que permita generar de forma automática los cuestionarios en otros idiomas.
	\item Integrar la posibilidad de crear cuestionarios desde cero desde la propia página web, contando con la ventaja de tener un modelo de datos ya definido.
	\item Se puede considerar la opción de migrar la aplicación a plataformas móviles, para no depender de la conexión a internet para acceder.
	\item Integrar un módulo de estadísticas, donde poder visualizar de forma gráfica los resultados obtenidos en los diferentes cuestionarios.
\end{itemize}





\bibliographystyle{plain}
\bibliography{bibliografia}

\end{document}
